\documentclass[10pt,a4paper]{amsart}
\usepackage[margin=19mm]{geometry}
\usepackage{amsmath,amssymb,mathtools}
\usepackage[scr=rsfs,cal=euler]{mathalfa}
\usepackage{xfrac}
\usepackage[unicode,hidelinks,bookmarks=false]{hyperref}
\newcommand{\ie}{i.e.\ }
\newcommand{\cf}{cf.\ }
\newcommand{\resp}{respectively\ }
\newcommand{\Subsection}{\subsection}
\providecommand{\nobreakdash}{\nobreak\hbox{-}}
\setlength{\emergencystretch}{2em}
\multlinegap0pt
\usepackage{amssymb}
\usepackage{enumerate}
\usepackage{float}
\usepackage{xcolor}
\usepackage{tikz}
\newtheorem{lem}{Lemma}
\newcommand{\bea}{\begin{eqnarray}}
\newcommand{\beas}{\begin{eqnarray*}}
\newcommand{\eea}{\end{eqnarray}}
\newcommand{\eeas}{\end{eqnarray*}}
\title{Sharp Estimates for Hankel, Fekete--Szeg\"o and Zalcman Functionals for $\mathcal{S}_\mathbb{B}^{*}(\alpha)$ in Complex Banach Spaces}
\author{Nabadwip Sarkar, Pradip Das}
\date{Source: arXiv:2607.17584v1 (CC BY 4.0)}
\begin{document}
\maketitle

\noindent\emph{Source and licence.} Sharp Estimates for Hankel, Fekete--Szeg\"o and Zalcman Functionals for $\mathcal{S}_\mathbb{B}^{*}(\alpha)$ in Complex Banach Spaces. Version arXiv:2607.17584v1 (CC BY 4.0), distributed by arXiv under the Creative Commons Attribution 4.0 International licence, http://creativecommons.org/licenses/by/4.0/, as recorded in the arXiv licence metadata for this exact version and shown on the abstract page https://arxiv.org/abs/2607.17584v1. Full licence notice: \texttt{licenses/arxiv-2607.17584-CC-BY-4.0.txt}.\par\smallskip

\noindent\textit{Editorial scope.} Counted unit: the article's lem L5 (source0.tex lines 423--429). The article's complete original proof of it is delivered with it (source0.tex lines 430--512, one \texttt{proof} environment of 83 lines). No mathematics is written, repaired or re-derived: the statement and the proof are the article's own lines, in the article's own notation, and no retained line is substituted or re-flowed.  Retained uncounted as the source-local context the proof needs: lines 74--78; lines 340--368.  Nothing was omitted from any retained span, so the package carries no omission record.  No retained line contains a citation, so the wrapper carries no bibliography and no bibliography entry.  No retained line contains a figure environment, an image-inclusion command or drawing code, so no image is delivered and the package declares no asset dependency.  Editorial formatting only, in addition: an AMS \texttt{amsart} preamble that restates the article's own declarations and macros used by the retained lines, namely the environments \texttt{lem} on separate counters, together with the standard packages those lines need.  The retained environments print in this document as Lemma 1 in order of appearance, whereas the article prints the same items with the numbering of its own full document.  The complete native source of the article as distributed in the arXiv e-print for this exact version is bundled unchanged as \texttt{sources/205581/source0.tex}.
\begin{equation}\label{eq:P}
p(z)=1+\sum_{m=1}^{\infty}c_mz^m
=1+c_1z+c_2z^2+c_3z^3+\cdots,
\qquad z\in\mathbb{U}.
\end{equation}

\begin{lem}\label{L4}  Let $g\in H(\mathbb{B},\mathbb{C})$ satisfy $g(0)=1$, and define $F(x)=g(x)x, x\in\mathbb{B}.$ Suppose that $F\in\mathcal{S}_{\mathbb{B}}^{*}(\alpha)$. Then, for every
$x_{0}\in X$ with $\|x_{0}\|=1$,
\bea\label{A234}
\begin{cases} A_{2}
&=
\frac{1}{2!}
T_{x_{0}}
\left(
D^{2}F(0)(x_{0}^{2})
\right)
=(1-\alpha)c_1,\\
A_{3}
&=
\frac{1}{3!}
T_{x_{0}}
\left(
D^{3}F(0)(x_{0}^{3})
\right)
=\frac{1}{2}(1-\alpha)\left(c_2+(1-\alpha)c_1^2\right),
\\
\text{and}\;\; 
A_{4}
&=
\frac{1}{4!}T_{x_{0}}\left(D^{4}F(0)(x_{0}^{4})\right)=\frac{1}{6}(1-\alpha)\left(2c_3+3(1-\alpha)c_1c_2+(1-\alpha)^2c_1^3\right),
\end{cases}
\eea
where $c_1$, $c_2$ and $c_3$ are given by (\ref{eq:P}).

\end{lem}

\begin{lem}\label{L5} Let $F$ be a locally biholomorphic mapping on $\mathbb{B}$ and $F$ satisfy the following assumptions:
\begin{equation}\label{cc1}
\frac{D^{k+1}F(0)\left(x^{k+1}\right)}{(k+1)!}
= H_{F,k}(x)x,\qquad x\in X,\quad k=1,2,3.
\end{equation}
where $H_{F,k}(x)$ is a homogeneous polynomial of degree $k$ with values in $\mathbb{C}$. If $F\in \mathcal{S}_{\mathbb{B}}^{*}(\alpha)$, then for each $x_0\in X$ with $\|x_0\|=1$,  the quantities $A_2$, $A_3$ and $A_4$ are given by  (\ref{A234}).
\end{lem}

\begin{proof} Fix $x_0\in \partial B$ and let $T_{x_0}\in T(x_0)$. Define
\bea\label{aaa1}
\phi(\xi)=
\begin{cases}
\dfrac{\xi}{T_{x_0}\!\left(\varphi(\xi x_0)\right)}, & \xi\neq 0,\\[2mm]
1, & \xi=0,
\end{cases}
\eea
where
\[
\varphi(x)=(DF(x))^{-1}F(x).
\]
Then $\phi\in H(\mathbb U)$ and $\phi(0)=1$. Moreover, since
$F\in \mathcal{S}_{\mathbb{B}}^{*}(\alpha)$ , it follows from the definition of Ozaki close to convex function that
\[
\Re\!\left(\frac{\xi}{T_{x_0}\!\left(\varphi(\xi x_0)\right)}\right)
=
\Re\!\left(\frac{\|\xi x_0\|}{T_{\xi x_0}\!\left(\varphi(\xi x_0)\right)}\right)>\alpha,\;\;\text{for}\;\;0\leq \alpha < 1
\qquad \xi\in\mathbb U\setminus\{0\}.
\]
Hence,
\[\Re\!\left(\phi(\xi)\right)>\alpha,\;\;\text{for}\;\;0\leq \alpha < 1\qquad \xi\in\mathbb U.
\]
 Then there exists $p\in \mathcal{P}$ such that 
\bea\label{kp1}\frac{1}{1-\alpha}(\phi(\xi)-\alpha)=p(\xi).\eea
Since $\phi$ is holomorphic function on $\mathbb{U}$, then it has Taylor series expansion:
\bea\label{phi1}\phi(\xi)=1+\frac{T_{x_0}(D^2\varphi(0)(x_0^2))}{2!}\xi+\frac{T_{x_0}(D^3\varphi(0)(x_0^3))}{3!}\xi^2+\frac{T_{x_0}(D^4\varphi(0)(x_0^4))}{4!}\xi^3+\cdots,\xi\in \mathbb{U}.\eea
Using (\ref{aaa1}), \eqref{eq:P}, \eqref{kp1} and \eqref{phi1}, a direct computation yields
\bea\label{qq1}
\begin{cases} \frac{T_{x_0}(D^2\varphi(0)(x_0^2))}{2!}&=-(1-\alpha)c_1,\\
 \frac{T_{x_0}(D^3\varphi(0)(x_0^3))}{3!}&=(1-\alpha)^2 c_1^2-(1-\alpha)c_2\\
\frac{T_{x_0}(D^4\varphi(0)(x_0^4))}{4!}&=-(1-\alpha)^3 c_1^3+2(1-\alpha)^2 c_1c_2-(1-\alpha)c_3.
\end{cases}
\eea
On the other hand, since $\varphi(x)=(DF(x))^{-1}F(x)$, it follows that
\bea\label{qqq2}
\begin{cases}
\frac{D^{2}F(0)(x^{2})}{2!}&=-\frac{D^{2}\varphi(0)(x^{2})}{2!}\\
 -\frac{D^{3}F(0)(x^{3})}{3}&=\frac{D^{3}\varphi(0)(x^{3})}{3!}+D^{2}F(0)\!\left(x,\,\frac{D^{2}\varphi(0)(x^{2})}{2!}\right)\\
-\frac{D^{4}F(0)(x^{4})}{8}&=\frac{D^{4}\varphi(0)(x^{4})}{4!}+D^{2}F(0)\!\left(x,\,\frac{D^{3}\varphi(0)(x^{3})}{3!}\right)\\
&\;\;\;+\frac{1}{2}D^{3}F(0)\!\left(x^{2},\,\frac{D^{2}\varphi(0)(x^{2})}{2!}\right).
\end{cases}\eea

Combining(\ref{qq1}) and (\ref{qqq2}) together with assumption \eqref{cc1}, we derive explicit formulas for the coefficients \(A_2\), \(A_3\)  and \(A_4\) in terms of the parameters \(p_1\), \(p_2\) and  \(p_3\), respectively.


\bea\label{A2}A_2&=&T_{x_0}\!\left(\frac{D^2F(0)(x_0^2)}{2!}\right)=H_{F,1}(x_0)
=-\,T_{x_0}\!\left(\frac{D^2\varphi(0)(x_0^2)}{2!}\right)
=(1-\alpha)c_1\eea
and
\[
\frac{D^2\varphi(0)(x_0^2)}{2!}=-\,\frac{D^2F(0)(x_0^2)}{2!}=-H_{F,1}(x_0)x_0
=-A_2x_0
=-(1-\alpha)c_1x_0.
\]

\bea\label{A3} A_3&=&T_{x_0}\!\left(\frac{D^3F(0)(x_0^3)}{3!}\right)
=H_{F,2}(x_0) \nonumber\\
&=&-\frac{1}{2}\left(T_{x_0}\!\left(\frac{D^3\varphi(0)(x_0^3)}{3!}\right)
+T_{x_0}\!\left(D^2F(0)\!\left(x_0,\frac{D^2\varphi(0)(x_0^2)}{2!}\right)\right)\right) \nonumber\\
&=&-\frac{1}{2}\left((1-\alpha)^2 c_1^2-(1-\alpha)c_2+T_{x_0}\!\left(D^2F(0)(x_0,-(1-\alpha)c_1x_0\right)\right) \nonumber\\
&=&-\frac{1}{2}\left((1-\alpha)^2 c_1^2-(1-\alpha)c_2-2(1-\alpha)c_1A_2\right)\nonumber\\
&=&\frac{1-\alpha}{2}\left(c_2+(1-\alpha) c_1^2\right)
\eea

and

\beas 
\frac{D^3\varphi(0)(x_0^3)}{3!}&=&-\frac{D^3F(0)(x_0^3)}{3}-D^2F(0)\!\left(x_0,\frac{D^2\varphi(0)(x_0^2)}{2!}\right)\\
&=&-\frac{D^3F(0)(x_0^3)}{3}-D^2F(0)(x_0,(1-\alpha) c_1 x_0)\\
&=&\left(-2H_{F,2}(x_0)-2(1-\alpha)^2 b_1^2\right) x_0\\
&=&(1-\alpha)((1-\alpha)c_1^2-c_2)x_0.\eeas
\bea\label{A4}
A_4&=& T_{x_0}\!\left(\frac{D^4F(0)(x_0^4)}{4!}\right)=H_{F,3}(x_0) \nonumber\\
&=&-\frac{1}{3}\Bigg(T_{x_0}\!\left(\frac{D^4\varphi(0)(x_0^4)}{4!}\right)+T_{x_0}\!\left(D^2F(0)\!\left(x_0,\frac{D^3\varphi(0)(x_0^3)}{3!}\right)\right) \nonumber\\
&& +\frac{1}{2}T_{x_0}\!\left(D^3F(0)\!\left(x_0^2,\frac{D^2\varphi(0)(x_0^2)}{2!}\right)\right)\Bigg) \nonumber\\
&=&-\frac{1}{3}\Bigg(-(1-\alpha)^3 c_1^3+2(1-\alpha)^2 c_1c_2-(1-\alpha)c_3 \nonumber\\
&&+
T_{x_0}\!\left(D^2F(0)(x_0,(1-\alpha)((1-\alpha)c_1^2-c_2)x_0)\right)+\frac{1}{2}T_{x_0}\!\left(D^3F(0)(x_0^2,-(1-\alpha)c_1x_0)\right)\Bigg) \nonumber\\
&=&-\frac{1}{3}\Bigg(-(1-\alpha)^3 c_1^3+2(1-\alpha)^2 c_1c_2-(1-\alpha)c_3\nonumber\\
&&+2(1-\alpha)((1-\alpha)c_1^2-c_2)x_0A_2-3(1-\alpha)c_1A_3\Big) \nonumber\\
&=&\frac{1}{6}(1-\alpha)\left((1-\alpha)^2 c_1 +3(1-\alpha)c_1c_2+2c_3\right).\eea
\end{proof}

\end{document}
